\section{Compression of distributed density networks}
\label{sec:compression}

We approximate the physical density produced by a distributed construction
after tracing out its private registers.  The private work spaces and
bipartite resources may have arbitrary finite dimensions.  We first replace
pair effects by operations using only prepared pair sources, then replace
the source operators in the density calculation by sums of product
operators.  The error is
controlled after contraction of the whole construction.  The result is an
operator tensor network, with one ket and one bra physical leg at each
party; it need not be positive.

We use a size parameter $L\ge2$.  A quantity is \emph{uniformly polynomial}
if it is at most $C L^a$, where $C,a$ are independent of the particular
construction and of all its private Hilbert-space dimensions.  These
constants may depend on a prescribed, fixed accuracy exponent.  Every
Hilbert space below is finite dimensional, but an unspecified dimension
has no size bound.

\subsection{Distributed constructions and pair effects}

A party is a label for a collection of registers.  At every time the
memory is a tensor product of the parties' current Hilbert spaces.  A
gate has specified input and output spaces for each participating party;
these spaces may differ.  Registers outside the participating parties
are spectators.  A \emph{private} gate has one participating party.
Gates are linear contractions.  Registers to be discarded are retained
as part of their owner's memory until the end, so the entire construction
is a pure linear calculation before its final partial trace.

An allowed monomial in a gate expansion is a finite composition of local
contractions, preparations of normalized vectors shared between two
participating parties, and contractions by normalized bras shared
between two participating parties.  A transferred register of dimension
$d$ can also be allowed when $d$ is uniformly polynomial: expand its
identity transfer as
\[
 \sum_{a=1}^{d}|a\rangle_{\mathrm{receiver}}
                 \langle a|_{\mathrm{sender}}.
\]
Each summand consists of local norm-one maps.  A bounded number of such
transfers therefore changes a polynomial gate expansion by only a
polynomial factor.  We make this expansion first whenever such messages
occur.

The density argument will sample prepared pair sources.  The following
lemma converts pair effects to that form while approximating the original
gate in operator norm.  The ideal extra output is the same fixed garbage
vector in every branch; this common comparison lets us sum the branchwise
operator-norm errors.

\begin{lemma}[Elimination of normalized pair effects]
\label{lem:effects}
Let $G$ be a contraction with an expansion
$G=\sum_{\xi}c_{\xi}M_{\xi}$ into $K$ allowed monomials on a fixed set of
parties.  Suppose each monomial has at most $r$ pair effects.  For every
integer $m\ge1$ there is a gate $\widetilde G_m$, with additional garbage
registers, and a fixed normalized garbage vector $\Gamma_m$, such that
\begin{equation}
 \norm{\widetilde G_m-G\otimes|\Gamma_m\rangle}
 \le \frac{r}{\sqrt m}\sum_{\xi}|c_{\xi}|.
 \label{eq:compression-effect-error}
\end{equation}
The gate $\widetilde G_m$ has an expansion using only local contractions
and normalized pair sources.  It has at most $K m^r$ monomials, and their
absolute coefficient sum is at most $\sum_{\xi}|c_{\xi}|$.  Its additional
registers are owned by the original participating parties.  All sources
on the same pair of parties may be combined into one normalized pair
source.
\end{lemma}

\begin{proof}
First consider a single normalized pair vector
$\eta\in\mathcal U\otimes\mathcal V$.  Its effect consumes a pair of
registers in this space.  Append $m-1$ copies of $\eta$, and average the
$m$ cyclic permutations of the $m$ pair registers.  Each permutation is
the tensor product of the same cyclic permutation of the $\mathcal U$
registers at one party and of the $\mathcal V$ registers at the other.
Denote the resulting map by
\[
 T_m:\mathcal U\otimes\mathcal V
       \longrightarrow(\mathcal U\otimes\mathcal V)^{\otimes m}.
\]
It is a contraction, being an average of isometries.  It sends $\eta$ to
$\eta^{\otimes m}$.  If $v\perp\eta$, its $m$ insertion positions among
the copies of $\eta$ are pairwise orthogonal and each has norm
$\norm{v}$.  Consequently
\begin{equation}
 \norm{T_m-|\eta\rangle^{\otimes m}\langle\eta|}
       =m^{-1/2}
 \label{eq:compression-one-effect}
\end{equation}
unless the perpendicular subspace is zero, in which case the left side
is zero.  This operator-norm estimate is unchanged by adjoining
spectators.  The permutation average has $m$ terms, each with coefficient
$1/m$ and consisting of local unitaries after normalized source
preparation.

For the whole gate, index every prescribed effect occurrence by its
monomial and its position within that monomial.  Give each such
occurrence its own output stack of $m$ pair registers, with the dimensions
and endpoint parties of its vector $\eta$.  In the branch using that
occurrence, replace its effect by $T_m$ and leave the stack idle
thereafter.  In every other branch append the full fixed stack
$\eta^{\otimes m}$ as a pair source.  Thus every branch has the same
additional output space, and its ideal additional output is the same
vector
\[
 \Gamma_m=\bigotimes_{\text{prescribed occurrences }o}
                    \eta_o^{\otimes m}.
\]
Different branches may consume different internal registers in their
original effects; the occurrence-specific output stacks still give a
common external output space.  No subsequent operation acts on a stack.

All original monomial factors, their ideal augmented replacements, and
the maps $T_m$ are contractions.  Telescoping the at most $r$ replacements
in one monomial bounds its error by $r/\sqrt m$.  Summing with the
original coefficients proves
\eqref{eq:compression-effect-error}.  Expanding the active permutation
averages creates at most $m^r$ terms per original monomial.  Their
averaging coefficients are nonnegative and sum to one, so this expansion
does not increase the absolute coefficient sum.  The inactive stacks
require no extra summation labels.  Finally, a tensor product of
normalized pair vectors on a fixed pair of parties is a normalized pair
vector on the tensor products of their respective half-spaces.
\end{proof}

\subsection{The density-compression theorem}

\begin{theorem}[Distributed density compression]
\label{thm:compression}
Consider a family of distributed constructions indexed by $L\ge2$ with
the following uniform properties.
\begin{enumerate}
\item The number of parties and gates is uniformly polynomial.  The
initial vector is a tensor product of normalized private vectors, one
at each party, and every gate is a contraction.
\item There is a constant $b$ such that each nonprivate gate involves at
most $b$ parties, and each party participates in at most $b$ nonprivate
gates over the entire construction.
\item Each nonprivate gate is a sum of uniformly polynomially many
allowed monomials with uniformly polynomially bounded scalar
coefficients.  There is a constant $r$ bounding the number of pair
sources, pair effects, and polynomial-dimensional messages in each
monomial.  All constituent local maps are contractions.
\item At the end, the intended physical output space of each party has
uniformly polynomial dimension.  All other registers are traced out
party by party.
\end{enumerate}
Write $\rho$ for the resulting physical output density operator; it is
positive and may have trace less than one.  For every fixed $p>0$, there
is an operator tensor network $\sigma$ satisfying
\[
 \norm{\sigma-\rho}_1\le L^{-p}.
\]
It has one node per party, a bounded number of incident virtual links at
each node, and uniformly polynomial link dimensions.  Every link joins
two parties which participate in a common nonprivate gate of the
original construction.  The physical legs at a node are the ket and bra
legs of that party's intended output.
The constants in these bounds depend only on $p$ and the fixed uniform
data in the assumptions, including $b,r$ and the polynomial bounds,
and not on private dimensions or source Schmidt ranks.

The same conclusion holds if the gate expansions in the third
assumption are available to every fixed inverse-polynomial
operator-norm accuracy, with uniform polynomial bounds at each fixed
accuracy.
\end{theorem}

\begin{proof}
We give the estimates for a target error $0<\epsilon\le1$ whose inverse
is uniformly polynomial; take $\epsilon=L^{-p}$ at the end.  Expand the
small messages as above.  This preserves all constant incidence and
participation bounds and gives a uniformly polynomial absolute
coefficient sum for every gate.

After eliminating effects, we will replace each source operator
$|\eta\rangle\langle\widetilde\eta|$ in the density calculation by an
unbiased sum of $k$ product operators.  We estimate the resulting physical
output, allowing the sampled factors themselves to have large norms.
Expanding the error by the positions where a centered correction is used
makes the lifetime participation bound decisive: $s$ corrected positions
have at most $2s$ endpoint parties, which together participate in at most
$2bs$ nonprivate gates.  We expand only those gates into monomials; all
other gates remain aggregate contractions.  Thus the coefficient and
physical-output choices cost a fixed polynomial per corrected position,
rather than a product over every gate in the construction.

For each fixed choice, the remaining task is a mean trace-norm estimate
of $k^{-s/2}$.  Normalized probability weights from the corrected endpoints
and the exact sources crossing the boundary of their party set make this
estimate
independent of all private dimensions.  Summing over correction positions
will then give a small error with polynomial $k$.

\paragraph{Reduction to source-only gates.}
Let $M$ bound the number of nonprivate gates.  If $M=0$, the output is a
product of private densities and the assertion is immediate.  Otherwise
choose $\delta=\epsilon/(8M)$.  Apply Lemma~\ref{lem:effects} at every
nonprivate gate, choosing one sufficiently large polynomial integer $m$
so that its error is at most $\delta$.  This choice is made using the
\emph{original} coefficient sum and number of effects, before expanding
the permutation averages.  Afterward there are polynomially many
monomials, and their absolute coefficient sum has not increased.  Thus
the choice of $m$ is not circular.

The gate $\widetilde G_m$ has norm at most $1+\delta$.  Replace it by
$\widetilde G_m/(1+\delta)$, a contraction within $2\delta$ of the
original gate with its fixed garbage appended.  Keep these garbage
registers idle.  The augmented original circuit has the same physical
output density as the original circuit, since all the appended garbage
vectors are fixed and normalized.  Telescoping products of contractions
shows that the two final pure vectors differ by at most $2M\delta$.
For vectors of norm at most one,
\[
 \norm{|v\rangle\langle v|-|w\rangle\langle w|}_1
       \le(\norm{v}+\norm{w})\norm{v-w}
       \le2\norm{v-w}.
\]
Partial trace does not increase this norm.  Hence the source-only
circuit's physical density, denoted $\rho'$, satisfies
\begin{equation}
 \norm{\rho'-\rho}_1\le4M\delta\le\epsilon/2.
 \label{eq:compression-effect-circuit-error}
\end{equation}
The number of gates and the list of participating parties have not
changed.  The new polynomial exponents are now fixed.  The sampling
parameter introduced below is chosen only after this reduction.

Each nonprivate gate of this new circuit can be written
\begin{equation}
 G_g=\sum_{\xi\in\Xi_g}c_{g,\xi}
       \Bigl(\bigotimes_{P}B_{g,P,\xi}\Bigr)
       \left[\Bigl(\bigotimes_{e\in E_g}
                         |\eta_{g,e,\xi}\rangle\Bigr)
                    \otimes\id_{\mathrm{in}}\right],
 \qquad \norm{G_g}\le1.
 \label{eq:compression-source-gate}
\end{equation}
Here $E_g$ contains one slot for each participating pair of parties,
each $\eta_{g,e,\xi}$ is normalized, and every $B_{g,P,\xi}$ is a local
contraction.  Unused pairs have a one-dimensional source.  To obtain
this form, prepare a monomial's sources in fresh registers before their
first use, and tensor the preceding local maps with identities on those
registers.  This preserves their norms and their chronological action.
Then combine sources on each pair, and compose the local maps at each
party.  This procedure stores fresh registers; it does not commute an
existing source half through an earlier map.

\paragraph{Common private slot spaces.}
The dimensions of a source slot can depend on $\xi$.  For each endpoint
$P$ of $e$, embed its half-space for branch $\xi$ isometrically into
\[
 \mathcal H_{g,P,e}
       =\bigoplus_{\xi\in\Xi_g}\mathcal H_{g,P,e,\xi}.
\]
Embed the branch source accordingly, and extend $B_{g,P,\xi}$ by
orthogonal projection onto its branch's source sectors.  Its norm
remains at most one.  All external gate input and output spaces are
already fixed by the memory layout and, in the preceding reduction, by
the common inventory of stacks.  Formula
\eqref{eq:compression-source-gate} therefore has one fixed ambient
space for each slot.  These direct sums are private spaces; no bound on
their dimensions will be used.

Fix uniformly polynomial bounds $B(L)$ and $d(L)$ such that
\begin{equation}
 \sum_{\xi\in\Xi_g}|c_{g,\xi}|\le B(L),
 \qquad
 \dim\mathcal H_{P,\mathrm{phys}}\le d(L)
 \label{eq:compression-polynomial-bounds}
\end{equation}
for all gates and parties.  The cardinalities $|\Xi_g|$ are also
uniformly polynomial.  We may increase $B,d$ to be at least one.

\paragraph{A random replacement for one density source.}
For a ket branch $\xi$ and bra branch $\zeta$ at a slot, write Schmidt
decompositions, in their respective bases, as
\[
 |\eta\rangle=\sum_a\sqrt{\lambda_a}
                      |u_a\rangle|v_a\rangle,
 \qquad
 |\widetilde\eta\rangle=\sum_c\sqrt{\mu_c}
                      |\widetilde u_c\rangle
                      |\widetilde v_c\rangle.
\]
Both probability lists sum to one; their lengths need not agree.  Let
$g_{ac}^{(j)}$, for $1\le j\le k$, be independent standard circular
complex Gaussians, normalized by $\mathbb E|g_{ac}^{(j)}|^2=1$.  Put
\begin{align}
 U_j&=\sum_{a,c}(\lambda_a\mu_c)^{1/4}g_{ac}^{(j)}
             |u_a\rangle\langle\widetilde u_c|,\nonumber\\
 V_j&=\sum_{b,d}(\lambda_b\mu_d)^{1/4}
                    \overline{g_{bd}^{(j)}}
             |v_b\rangle\langle\widetilde v_d|,\nonumber\\
 \widehat X&=\frac1k\sum_{j=1}^k U_j\otimes V_j,
 \qquad X=|\eta\rangle\langle\widetilde\eta|,
 \qquad \Delta=\widehat X-X.
 \label{eq:compression-random-source}
\end{align}
These matrices act between the corresponding bra and ket source
supports in the common ambient half-spaces.  No bound on their sampled
operator norms is claimed or needed.  Since
$\mathbb E[g_{ac}\overline{g_{bd}}]=\delta_{ab}\delta_{cd}$,
$\mathbb E\widehat X=X$.

The centered coefficient of $\Delta$ with ket endpoint indices $(a,b)$
and bra endpoint indices $(c,d)$ is
\[
 (\lambda_a\lambda_b\mu_c\mu_d)^{1/4}
 \left(\frac1k\sum_{j=1}^k
          g_{ac}^{(j)}\overline{g_{bd}^{(j)}}
                   -\delta_{ab}\delta_{cd}\right).
\]
For composite Gaussian indices $\alpha,\beta,\gamma,\delta$, the
complex Gaussian pairing formula \cite{Isserlis1918,FassinoPistoneRogantin2019}
gives the fourth-moment identity
\begin{equation}
 \mathbb E\left[
  (g_\alpha\overline{g_\beta}-\delta_{\alpha\beta})
  \overline{(g_\gamma\overline{g_\delta}
                         -\delta_{\gamma\delta})}\right]
       =\delta_{\alpha\gamma}\delta_{\beta\delta}.
 \label{eq:compression-gaussian-covariance}
\end{equation}
Indeed independence, circular symmetry, and $\mathbb E|g_\alpha|^4=2$
give the two Gaussian pairings
$\delta_{\alpha\beta}\delta_{\gamma\delta}
 +\delta_{\alpha\gamma}\delta_{\beta\delta}$, and centering removes
the first.  Thus the coefficient covariance is diagonal and its
diagonal entry is
\begin{equation}
 \frac1k\sqrt{\lambda_a\lambda_b\mu_c\mu_d}.
 \label{eq:compression-slot-variance}
\end{equation}
Take fresh independent samples for each slot location and each ordered
pair of monomial labels at that gate.

\paragraph{Expansion by corrected positions.}
Let $\mathcal E=\{(g,e):e\in E_g\}$ be the set of slot locations and
$N=|\mathcal E|$.  It is uniformly polynomial.  In the density
calculation for each gate, substitute $X+\Delta$ for each source
operator inside its sum over ket and bra labels.  Expanding these
products throughout the circuit gives
\begin{equation}
 \widehat\rho-\rho'
       =\sum_{\varnothing\ne S\subseteq\mathcal E}F_S.
 \label{eq:compression-subset-expansion}
\end{equation}
A set $S$ specifies which \emph{positions} use centered corrections;
it does not specify their monomial labels.  At a gate with no selected
position, summing its labels gives exactly $G_g(.)G_g^\dagger$.

Fix $S$, write $s=|S|$, and let $\mathcal A$ be the parties which are
endpoints of its slots.  Then $|\mathcal A|\le2s$.  Expand the ket and
bra labels at every nonprivate gate touching $\mathcal A$, whether or
not that gate has a selected slot.  There are at most $2bs$ such gates.
All other gates lie entirely in the exterior and remain their exact
aggregate contractions.  Also fix ket and bra basis entries $x,y$ of
the physical outputs at parties in $\mathcal A$.  The total weighted
cost of these choices is at most
\begin{equation}
 B(L)^{4bs}d(L)^{4s}=Q(L)^s,
 \qquad Q(L):=B(L)^{4b}d(L)^4.
 \label{eq:compression-choice-cost}
\end{equation}
Here each expanded gate contributes at most
$(\sum_\xi|c_{g,\xi}|)^2$; the number of physical matrix entries is at
most $d(L)^{2|\mathcal A|}$.  Scalar coefficients are removed from the
maps and charged in this bound.

From now on fix one such deterministic choice of labels and physical
entries.  The samples used at its different corrected positions are
independent.  There is no sample-dependent selection of a branch.  We
will apply the triangle inequality over these choices, so no
cross-branch covariance assumption is necessary.  It remains to bound the
mean trace norm for this fixed choice by $k^{-s/2}$, with no private-dimension
factor; the normalized source weights will provide that bound.

\paragraph{The two separated contractions.}
Both ends of every corrected source lie in $\mathcal A$.  All sources
crossing from $\mathcal A$ to its exterior are therefore exact.  Their
tensor product is a normalized pure vector on the ket side and a
possibly different normalized pure vector on the bra side.  In their
Schmidt bases let the probability matrices be $\tau$ and
$\widetilde\tau$.  Their diagonal entries sum to one.  Empty lists of
sources are represented by a one-dimensional space and the scalar
probability one.

For every corrected slot retain both ket endpoint indices as free
indices; collect them in $i$.  Similarly collect the bra endpoint
indices in $l$.  Define normalized product probability lists $p_i$ and
$\widetilde p_l$ by giving a corrected slot the respective factors
$\lambda_a\lambda_b$ and $\mu_c\mu_d$.  In particular the two endpoint
indices are independent summation indices, even though an exact pure
source would identify them.  If $z_{il}$ is the product of the centered
coefficients at these slots, independence and
\eqref{eq:compression-slot-variance} give
\begin{equation}
 \mathbb E[z_{il}\overline{z_{i'l'}}]
       =k^{-s}\sqrt{p_i\widetilde p_l}\,
                      \delta_{ii'}\delta_{ll'}.
 \label{eq:compression-product-covariance}
\end{equation}

Regard all corrected endpoint registers and all affected-side halves of
exact crossing sources as free inputs supplied at the beginning of the
separated calculation.  This again means keeping fresh registers idle
until first use.  Prepare all exact sources internal to either side as
normalized vectors on that side.  Every previously expanded crossing
gate now consists of local contractions and these exact sources.  Thus
there is no remaining communication between the two sides.  Gates
wholly in the exterior which were not expanded remain their aggregate
contractions; their internal branch sums are not opened.

Let $\mathcal I$ be the ket space with basis $i$ and $\mathcal T$ the
affected-side ket Schmidt-support space with basis $t$ for $\tau$.
Use tildes for their bra analogues.  Including final physical basis
projection, the affected calculations give contractions
\[
 A_x:\mathcal I\otimes\mathcal T\longrightarrow\mathcal F,
 \qquad
 \widetilde A_y:\widetilde{\mathcal I}\otimes
              \widetilde{\mathcal T}\longrightarrow\mathcal F.
\]
The codomain $\mathcal F$ contains all affected discard registers.  It
is the same on ket and bra because the external memory spaces are fixed;
the inventory construction ensured this also for the added garbage.
These are contractions on \emph{all} free inputs: each operation in
their construction is a contraction, an isometric preparation, an
identity on an idle register, or a physical basis projection.  No
normalization of an output vector has occurred.  Therefore
\begin{equation}
 O:=\widetilde A_y^\dagger A_x:
       \mathcal I\otimes\mathcal T\longrightarrow
       \widetilde{\mathcal I}\otimes\widetilde{\mathcal T},
 \qquad \norm{O}\le1.
 \label{eq:compression-block-contraction}
\end{equation}
Write its blocks as $O_{li}:\mathcal T\to\widetilde{\mathcal T}$.

The exterior calculations similarly give contractions from their
crossing-source half-spaces to the exterior physical output and
garbage.  In the matching Schmidt coordinate spaces denote them by
$K$ and $\widetilde K$.  They can differ on ket and bra but do not
depend on $i,l$.  Those indices are free inputs only on the affected
side, and fixing all crossing-gate labels removed every other means
of communication across the division.  For any rectangular operator
$Z$ between the two exterior input spaces,
\begin{equation}
 \norm{\Tr_{\mathrm{garbage}}(KZ\widetilde K^\dagger)}_1
          \le\norm{Z}_1.
 \label{eq:compression-exterior-contraction}
\end{equation}
To see this, test the output against an arbitrary matrix $M$ of operator
norm at most one.  The pulled-back test is
$\widetilde K^\dagger(M\otimes\id)K$, of norm at most one.
Trace-norm duality proves the inequality even when $Z$ is neither
Hermitian nor positive.

\paragraph{The second-moment bound.}
Density-weighted comparisons between trace and Hilbert--Schmidt norms
also occur in privacy amplification and decoupling
\cite[Lemma~5.1.3]{Renner2005Security} \cite[Lemma~3.7]{DupuisBertaWullschlegerRenner2014}.
Here we use a two-weight rectangular form of Schatten H\"older, proved below.
Tracing the affected garbage gives the following exterior input
operator for the fixed labels and $x,y$:
\begin{equation}
 Z=\sum_{i,l}z_{il}\,
          \tau^{1/2}O_{li}^{T}\widetilde\tau^{1/2}.
 \label{eq:compression-exterior-input}
\end{equation}
Indeed its $(t,t')$ coefficient is
$\sum_{i,l}z_{il}\sqrt{\tau_t\widetilde\tau_{t'}}
       \langle\widetilde A_y(l,t'),A_x(i,t)\rangle$.
The transpose in \eqref{eq:compression-exterior-input} is the ordinary
full matrix transpose in these bases, including when the matrix is
rectangular; it preserves all Schatten norms.

For an operator $W$ between these coordinate spaces, Schatten
H\"older with exponents $4,2,4$ gives
\begin{equation}
 \norm{\tau^{1/2}W\widetilde\tau^{1/2}}_1
 \le\norm{\tau^{1/4}}_4\,
       \norm{\tau^{1/4}W\widetilde\tau^{1/4}}_2\,
       \norm{\widetilde\tau^{1/4}}_4
 =\norm{\tau^{1/4}W\widetilde\tau^{1/4}}_2.
 \label{eq:compression-quarter-powers}
\end{equation}
The subscript $2$ here denotes the Hilbert--Schmidt norm, and the
quarter-power Schatten-$4$ norms equal one because the probability
matrices have trace one.  Applying this and transposing the middle
matrix in \eqref{eq:compression-exterior-input} yields
\[
 \norm{Z}_1\le
 \norm{\sum_{i,l}z_{il}
             \widetilde\tau^{1/4}O_{li}\tau^{1/4}}_2.
\]
By \eqref{eq:compression-product-covariance}, the expected square of
the right side is
\begin{equation}
 k^{-s}\sum_{i,l}\sqrt{p_i\widetilde p_l}\,
          \norm{\widetilde\tau^{1/4}O_{li}\tau^{1/4}}_2^2.
 \label{eq:compression-block-second-moment}
\end{equation}
Define trace-one positive diagonal matrices on the full free input
spaces by
\[
 R=\operatorname{diag}(p_i)\otimes\tau,
 \qquad
 \widetilde R=\operatorname{diag}(\widetilde p_l)
                                      \otimes\widetilde\tau.
\]
The sum in \eqref{eq:compression-block-second-moment} is exactly
$\norm{\widetilde R^{1/4}OR^{1/4}}_2^2$: its $(l,i)$ block is
$\widetilde p_l^{1/4}p_i^{1/4}
 \widetilde\tau^{1/4}O_{li}\tau^{1/4}$, and the
Hilbert--Schmidt squares of disjoint blocks add.  Another application
of Schatten H\"older and \eqref{eq:compression-block-contraction} gives
\begin{equation}
 \norm{\widetilde R^{1/4}OR^{1/4}}_2
   \le\norm{\widetilde R^{1/4}}_4\norm{O}\norm{R^{1/4}}_4
   \le1.
 \label{eq:compression-dimension-free}
\end{equation}
This bound contains no rank or dimension of a source, an input, or a
discard space.  Cauchy--Schwarz for the expectation now proves
\begin{equation}
 \mathbb E\norm{Z}_1\le k^{-s/2}.
 \label{eq:compression-one-choice}
\end{equation}
Applying the exterior map preserves this bound by
\eqref{eq:compression-exterior-contraction}.  Reinserting the affected
physical matrix unit $|x\rangle\langle y|$ preserves the trace norm,
since that matrix unit has trace norm one.

\paragraph{Summation and realization as a network.}
Summing over the fixed choices with
\eqref{eq:compression-choice-cost} gives
\[
 \mathbb E\norm{F_S}_1\le Q(L)^{|S|}k^{-|S|/2}.
\]
Consequently
\begin{equation}
 \mathbb E\norm{\widehat\rho-\rho'}_1
 \le\left(1+\frac{Q(L)}{\sqrt k}\right)^N-1.
 \label{eq:compression-total-error}
\end{equation}
If $N=0$, there is no sampling error.  Otherwise choose an integer
$k$ with
\[
 \frac{Q(L)}{\sqrt k}\le\frac{\epsilon}{8N}.
\]
It is uniformly polynomial.  The right side of
\eqref{eq:compression-total-error} is at most
$e^{\epsilon/8}-1\le\epsilon/4$.  Hence there is a realization with
$\norm{\widehat\rho-\rho'}_1\le\epsilon/2$, for example by Markov's
inequality.  Fix one such realization and set $\sigma=\widehat\rho$.
Together with \eqref{eq:compression-effect-circuit-error}, this proves
the asserted error.

It remains to check the network dimensions and link locality.  At a
gate, the common label $(\xi,\zeta)$ ranges over a polynomial set.
Choose one participating party and share this label along a star to
the other participating parties, using one virtual index of that
polynomial dimension on each star edge.  Each slot's sample index
$j\in\{1,\ldots,k\}$ is shared directly between its two endpoint
parties.  The local source matrices in
\eqref{eq:compression-random-source} are determined by these labels
and the fixed samples.  Put the gate's scalar coefficient and averaging
factors into any one local tensor.  All local maps, private source
preparations, intermediate memories, and final garbage traces can then
be contracted within each party to give a single tensor with its
physical ket and bra legs.  In particular, coordinates in the common
private direct sums are contracted locally and are not virtual link
indices.

There are at most $b-1$ star links and $\binom b2$ sample links per
gate.  Every party belongs to at most $b$ nonprivate gates, so the
resulting network has bounded incidence.  Each link joins two
participants of one original gate.  Parallel links can be left
separate or combined; only boundedly many meet a party, so combining
them preserves a polynomial dimension bound.  Sampled tensor entries
need not be bounded.  Their existence and link dimensions, rather than
their numerical size or efficient computation, are what is required.

Finally, suppose only approximate gate expansions are provided.  At
any prescribed fixed accuracy exponent their operator-norm errors can
be made smaller than $\epsilon$ divided by a sufficiently large
polynomial bound on the number of gates.  A near-contraction can be
rescaled by its known bound $1+\delta$ to be a contraction at an
additional error at most $\delta$.  The same telescoping estimate as
above absorbs these errors.  Apply the proved result at a fixed
smaller error budget.  All subsequent choices remain uniform
polynomials at that precision, which proves the last assertion.
\end{proof}

\begin{remark}[Absolute error and retained norm]
\label{rem:compression-retained-norm}
Theorem~\ref{thm:compression} approximates the subnormalized density
produced by the contraction circuit.  It does not normalize that
density.  In its later application, the construction separately ensures
that its physical output is close in trace norm to a normalized pure
state; in particular, the retained norm is close to one.  The theorem's
absolute error can then be added directly to that approximation error.
No bound on the private dimensions or on the Schmidt ranks of the
supplied pair vectors is needed for this use.
\end{remark}
